\documentclass[10pt,a4paper,twocolumn]{article}
\usepackage[utf8]{inputenc}
\usepackage[T1]{fontenc}
\usepackage[english]{babel}
\usepackage[top=2cm,bottom=2cm,left=1.8cm,right=1.8cm]{geometry}
\usepackage{booktabs}
\usepackage{graphicx}
\usepackage{hyperref}
\usepackage{natbib}
\usepackage{setspace}
\setlength{\columnsep}{0.6cm}

\title{\textbf{Parcel Lockers and Pickup Points in Urban Logistics:\\
Location Criteria and Delivery Scheme Design\\
--- A Case Study of La Rochelle, France}}

\author{[Author Names]\\
\small EIGSI La Rochelle, France\\
\small \texttt{[email@eigsi.fr]}}

\date{}

\begin{document}
\maketitle
\thispagestyle{empty}

%% ─────────────────────────────────────────────
\section*{Abstract}
The rapid growth of e-commerce multiplies urban delivery addresses and failed
delivery attempts. Parcel lockers and pickup points (PPPs) promise to consolidate
remises and reduce last-kilometre costs. However, their net benefit depends on
network design, upstream supply chain integration, and recipient access mode. This
paper proposes a multicriteria location framework and an integrated delivery scheme
(micro-hub + cargo cycles + PPPs) grounded in the scientific literature and applied
to La Rochelle, France, using territorial diagnoses of 2018 and 2024. Key findings
show that site performance is conditional on integration into existing travel chains,
upstream consolidation via a micro-hub, and cross-operator sharing of infrastructure.

\textbf{Keywords:} parcel lockers; pickup points; urban logistics; last-mile
delivery; micro-hub; location optimisation; La Rochelle.

%% ─────────────────────────────────────────────
\section{Introduction}

Home delivery failure rates of 10–30\% generate redundant trips and emissions.
Parcel lockers and pickup points (PPPs) address this by consolidating deliveries at
shared access points. Their net impact, however, depends on where they are placed and
how they are supplied. This paper asks: (1) What criteria determine PPP location
performance? (2) How should an integrated delivery scheme be designed for a
medium-sized French city?

We answer by synthesising the scientific literature and applying it to La Rochelle
(Charente-Maritime, France, $\approx$80,000 inhabitants), drawing on flow data from
a 2018 FRETURB-based freight diagnosis and spatial recommendations from a 2024
logistics land study, both produced by Interface Transport for the Community of
Agglomeration (CdA) of La Rochelle.

%% ─────────────────────────────────────────────
\section{State of the Art}

\subsection{Consolidation Benefits}

\citet{morganti2014} showed that relay points and lockers developed in France and
Germany to reduce uncertainty from recipient absence. \citet{iwan2016} demonstrated
density and productivity gains on the Polish InPost network. \citet{seghezzi2022}
modelled lower delivery costs versus home delivery in urban contexts, driven by
increased stop density.

\subsection{Location as a Network Problem}

\citet{deutsch2018} formalised the joint optimisation of number, location, and size
of lockers. \citet{zhou2024} differentiated home-based users from commuters, showing
that sites on existing travel chains generate less additional car traffic. \citet{castillo2024}
demonstrated in Barcelona that micro-hub location and routing must be co-optimised
for kilometre reduction.

\subsection{User Accessibility}

Acceptability depends on proximity, hours, security, and cycling/walking access
\citep{kedia2017}. \citet{vakulenko2018} note that lockers transfer logistics work
to the consumer. \citet{mehmood2022} found that $>$71\% of Nanjing locker sites
were primarily car-accessible, flagging the limits of distance-only analysis.

\subsection{Conditional Environmental Benefit}

\citet{schnieder2021} show that the emission balance is sensitive to recipient travel
mode. \citet{gutenschwager2024} stress adding operator tours and retrieval trips.
\citet{silva2023}, reviewing 102 papers, confirm benefits in distance and failed
deliveries but highlight dedicated car trips as a recurrent risk. \citet{ghazal2025},
modelling Aachen, found that a combined scenario (electric vehicles + lockers +
relays) reduced costs by $\geq$5.4\% and CO$_2$ by $\geq$61.1\%. \citet{pinchasik2025}
corroborate cost, traffic, and emission reductions for expanded locker use in the
Greater Oslo area.

\subsection{Micro-Hubs and Cargo Cycles}

\citet{katsela2022}, analysing 17 European and North American cases, show that
micro-hubs are a family of transhipment interfaces whose success depends on
stakeholder organisation and local context. They enable cargo-cycle distribution in
dense zones, provided sufficient volume and cross-operator cooperation are secured.

%% ─────────────────────────────────────────────
\section{Case Study: La Rochelle}

\subsection{Territorial Context}

La Rochelle is a mid-sized coastal city with a historic centre, a summer tourism
peak, and a CdA of $\approx$200,000 inhabitants. The city has invested in cycling
infrastructure and aims to reduce urban freight externalities.

\subsection{2018 Freight Diagnosis}

The FRETURB model estimates \textbf{101,325 weekly B2B movements} across the CdA,
of which \textbf{49,898} concern La Rochelle and \textbf{12,126} the city centre.
E-commerce adds \textbf{25,300 weekly deliveries} at CdA level, $\approx$\textbf{2,000
per week in the city centre}. Approximately \textbf{950 freight vehicles} access the
city centre daily, covering $\approx$\textbf{8,640 km/day} (780 vehicles
$<$3.5\,t; 170 heavy vehicles). Two demand poles emerge: the dense
\textit{Centre~Marché} sector and the collective housing zone of
\textit{Mireuil-Europe} \citep{interfacetransport2018}.

\subsection{2024 Logistics Land Study}

The 2024 study translates flows into spatial needs to 2040. It proposes a
hierarchical network: peripheral urban distribution platforms, urban distribution
spaces (EUD, 150\,m$^2$ to several thousand m$^2$) in the city core, and
cycle-logistics hubs within 2\,km of the centre (150–800\,m$^2$, ground floor,
cycle access, battery charging). Priority investigation sites include the Quartier
Gare/former SERNAM site and several underutilised industrial premises
\citep{interfacetransport2024}.

%% ─────────────────────────────────────────────
\section{Proposed Framework}

\subsection{Multicriteria Location Criteria}

Drawing on the literature and territorial data, we propose nine weighted criteria
(Table~\ref{tab:criteria}). An eliminating pre-screen removes sites without
accessible pedestrian routes, safe unloading, or adequate space.

\begin{table}[h]
\centering
\caption{Multicriteria site selection framework}
\label{tab:criteria}
\small
\begin{tabular}{lc}
\toprule
\textbf{Criterion} & \textbf{Weight} \\
\midrule
B2C/B2B demand potential & 20\% \\
Walking and cycling accessibility & 15\% \\
Integration into existing travel chains & 15\% \\
Operator supply access & 15\% \\
Network complementarity (gap coverage) & 10\% \\
Urban integration, safety, accessibility & 10\% \\
Technical feasibility & 5\% \\
Land and economic feasibility & 5\% \\
Territorial and social equity & 5\% \\
\bottomrule
\end{tabular}
\end{table}

\subsection{Site Typology for La Rochelle}

Five site types address different demand segments:
\textit{mobility lockers} (station, bus hubs, park-and-ride);
\textit{residential lockers} (Mireuil-Europe dense housing);
\textit{commercial lockers} (supermarkets, markets);
\textit{relay points} (proximity shops, human service, returns);
\textit{workplace lockers} (major employers, campus).

\subsection{Integrated Delivery Scheme}

The proposed scheme operates as a hierarchical chain:
(1) inbound flows are consolidated at a peripheral platform;
(2) city-bound goods transfer to a micro-hub or EUD near the centre
(Quartier~Gare/SERNAM);
(3) the micro-hub sorts by channel and dispatches via cargo cycles or electric
light vehicles;
(4) lockers and relay points double as return collection points, enabling
reverse logistics on outbound tours.

This avoids the key failure mode identified in the literature: deploying lockers
without upstream consolidation, which produces dedicated carrier runs to individual
units and cancels density gains.

%% ─────────────────────────────────────────────
\section{Conclusion}

The impact of PPPs depends less on the technology than on three design choices:
\textit{where} to locate them (sites that capture existing travel chains);
\textit{how} to supply them (consolidated from a micro-hub, not by individual
vans); and \textit{how} to govern them (shared infrastructure, interoperable
systems, public monitoring).

For La Rochelle, the 2018 and 2024 data quantify the problem (950 vehicles/day,
8,640~km/day in the city centre) and identify candidate sites. The proposed
framework provides a replicable decision-support method applicable to comparable
European mid-sized cities.

A reversible pilot combining one micro-hub, two to three locker sites, and
instrumented cargo-cycle tours is the recommended next step. Evaluation metrics
should cover operator kilometres, user modal split, fill rates, and full CO$_2$
balance including retrieval trips.

%% ─────────────────────────────────────────────
\bibliographystyle{apalike}
\bibliography{city_logistics_lockers_refs}

\end{document}
